% !PRINTED 181
\weberpage{181}

wir aber diese Form nach Reihe (VIII) \S~7, so kommt:

\weq{1}{(b\Delta)^3 b_y \Delta_x = [(b\Delta)^3 b_x \Delta_x]_y + \frac{1}{2}(b\Delta)^4(yx),}

oder weil:
\[
\begin{gathered}
(b\Delta)^3 b_x \Delta_x = (f,\Delta)^3 = 0 \\
(b\Delta)^4 = (f,\Delta)^4 = j,
\end{gathered}
\]
so ergiebt sich nach Ausführung der erwähnten Substitution

\weq{2}{(ab)(b\Delta)^3 \Delta_x a_x^3 = \frac{1}{2} j\cdot f,}

und demnach durch Comparation von (1) und (2)

\weq{3}{(\Delta,\Delta)^2 = \frac{1}{3}\cdot j\cdot f - \frac{i}{6}\Delta.}

159. \textit{Berechnung von $(\Delta,\Delta)^4$.} Der Werth dieser Invariante wird am einfachsten erhalten, wenn wir in der Ueberschiebung
\[
(f,\Delta)^2 = \frac{i}{6}\cdot a_x^4 = (a\Delta)^2 a_x^2 \Delta_x^2
\]
$x$ durch $b$, und gleichzeitig in der Identität
\[
\Delta_x^4 = (ab)^2 a_x^2 b_x^2
\]
$x$ durch $\Delta$ ersetzen. Wir erhalten alsdann:
\[
\begin{gathered}
\frac{i}{6}\cdot i = (a\Delta)^2(ab)^2(\Delta b)^2 \\
(\Delta,\Delta)^4 = (ab)^2(a\Delta)^2(b\Delta)^2,
\end{gathered}
\]
und demnach durch Comparation:

\weq{4}{(\Delta,\Delta)^4 = \frac{i^2}{6}.}

Hiemit ist der erste Theil unserer Aufgabe erledigt. Es zeigte sich, dass in der That die Ueberschiebungen $(f,\Delta)^\lambda$ und $(\Delta,\Delta)^\lambda$ keine neuen Formen erzeugen. Im zweiten Theile haben wir nunmehr nachzuweisen, dass
\[
(ab),\ (a\Delta)^2,\ (\Delta\Delta_1),\ (a\Delta)(b\Delta)
\]
Reducenten sind.

160. \textit{Die Reducenten.} Auf Grund der nunmehr berechneten Ueberschiebungswerthe ist nicht schwer nachzuweisen, dass zunächst $(ab)$, $(a\Delta)^2$, $(\Delta\Delta_1)$ Reducenteneigenschaft besitzen. Denn entwickelt man die entsprechenden einfachsten symbolischen Producte
\[
(ab)a_x^3 b_y^3,\ (a\Delta)^2 a_x^2 \Delta_y^2,\ (\Delta\Delta_1)\Delta_x^3\Delta_{1y}^3
\]
in die Reihe (VIII) \S~7, so findet man, dass die rechts auftretenden Elementarcovarianten sämmtlich Formen des Systems sind. Man erhält nämlich:
%CONT>>


% ===== pdf idx 196, printed 182.0 =====
% !PRINTED 182
\weberpage{182}

%<<CONT
\weq{1}{(ab)a_x^3 b_y^3 = \frac{3}{2}\Delta_{y^2}(xy) + \frac{1}{4} i(xy)^3}

\weq{2}{(a\Delta)^2 a_x^2 \Delta_y^2 = \frac{i}{6} f_{y^2} + \frac{1}{3} j\cdot(xy)^2}

\weq{3}{(\Delta\Delta_1)\Delta_x^3\Delta_{1y}^3 = \frac{1}{2}\left(j\cdot f_{y^2} - \frac{i}{2}\Delta_{y^2}\right)(xy) + \frac{1}{24} i^2 (xy)^3.}

Nun kann jedes symbolische Product, das $(ab)$, oder $(a\Delta)^2$, oder $(\Delta\Delta_1)$ zum Factor hat, durch Ueberschiebung mit den Formen (1), resp. (2), (3) oder deren Polaren gewonnen werden (vergl. \S~3) und die Gleichungen (1), (2), (3) lehren daher:

1) Alle Formen, welche $(ab)$ zum Factor haben, lassen sich reduciren auf Formen, welche an Stelle der beiden Symbole $a$ und $b$ nur mehr das eine $\Delta$ besitzen und auf Formen, die $i$ zum Factor haben.

2) Alle Formen, welche $(a\Delta)^2$, $(\Delta\Delta_1)$ besitzen, lassen sich auf Formen reduciren, welche $i$ oder $j$ zum Factor haben,

und daher sind die erwähnten Klammerfactoren in der That Reducenten. (Vergl. Nr. 124.)

161. Um nun die gleiche Eigenschaft auch für $(a\Delta)(b\Delta)$ nachzuweisen, denken wir uns wieder das einfachste symbolische Product, das $(a\Delta)(b\Delta)$ zu Klammerfactoren besitzt, nämlich
\[
P = (a\Delta)(b\Delta)a_x^3 b_y^3 \Delta_z^2
\]
in eine nach Potenzen von $(xy)$ fortschreitende Reihe entwickelt. Dann enthalten alle Elementarcovarianten, abgesehen von der ersten, den Factor $(ab)$ und sind demnach reducibel. Nun ist aber nach dem Productsatz diese erste Elementarcovariante:
\[
(a\Delta)(b\Delta)a_x^3 b_x^3 = \frac{1}{2} a_x^2 b_x^2 \{(b\Delta)^2 a_x^2 + (a\Delta)^2 b_x^2 - \Delta_x^2(ab)^2\},
\]
und somit, weil rechts nur Reducenten sich befinden, gleichfalls reducibel. Es ist also auch $P$ reducibel und $(a\Delta)(b\Delta)$ Reducent. Damit sind wir aber am Schlusse unseres Beweises angelangt. Jedes symbolische Product, das eine Co- oder Invariante der biquadratischen Form darstellen soll, kann zunächst reducirt werden auf eine Summe von Producten, die $i$ oder $j$ zum Factor haben, oder wenigstens auf Producte, mit weniger Symbolen $a$ oder $\Delta$. Durch Wiederholung der Reduction wird schliesslich jedes symbolische Product reducirt auf eine rationale ganze Function der Formen $f$, $\Delta$, $t$, $i$, $j$.

162. \textit{Die Fundamentalrelation zwischen den fünf Formen des Systemes.} Zwischen den fünf Formen des Systemes besteht eine Relation, die
%CONT>>


% ===== pdf idx 197, printed 183.0 =====
% !PRINTED 183
\weberpage{183}

%<<CONT
wir sofort nach allgemeinen Gesetzen berechnen können. Die Covariante $t$ sechsten Grades ist nämlich Functionaldeterminante von $f$ und $\Delta$, und demnach ist ihr Quadrat durch die geraden Formen des Systemes darstellbar. Man erhält nach Nr. 48, (I):
\[
2t^2 = \left|
\begin{array}{ccc}
\Delta, & \dfrac{i}{6} f, & f \\[4pt]
\dfrac{i}{6} f, & \dfrac{j}{6} f - \dfrac{i}{6}\Delta, & \Delta \\[4pt]
f, & \Delta, & 0
\end{array}
\right|
\]
oder:
\[
2t^2 = \frac{i}{2} f^2 \cdot \Delta - \Delta^3 - \frac{j}{3} f^3.
\]

163. \textit{Relationen zur Berechnung der Ueberschiebung von $t$ über sich selbst und über $f$ und $\Delta$.} Aus den letzten Untersuchungen geht hervor, dass die Ueberschiebungen von $t$ über sich selbst und über $f$ und $\Delta$ zu keinen neuen Covarianten Veranlassung geben können. Wenn wir im Folgenden sie dennoch berechnen, so geschieht dies einmal, weil sich uns hier ein lehrreiches Beispiel für die Behandlung derartiger Aufgaben darbietet, dann aber auch, weil gerade die Ueberschiebungen von $t$ über sich selbst interessante Aufschlüsse über den Charakter dieser Covariante gewähren. Am besten wählt man zur Bildung der durch diese Ueberschiebungen entstehenden Formen als Ausgangspunkt die Form:
\[
P = a_y^4 \Delta_x^4 - a_x^4 \Delta_y^4 = p_y^4 q_x^4,
\]
welche wir schon früher als Combinante charakterisirt haben. Entwickelt man dieselbe nach Potenzen von $(yx)$, so erhält man

\weq{1}{P = \sum \frac{\binom{4}{\varrho}\binom{4}{\varrho}}{\binom{8-\varrho+1}{\varrho}} \cdot \overline{(pq)}^{\,\varrho}_{y^{4-\varrho}} (yx)^\varrho.}

Von den Elementarcovarianten dieser Form $P$ verschwinden aber alle bis auf die zweite:
\[
(\overline{pq})^1 = [(a\Delta)a_x^3 \Delta_x^3 - (\Delta a)a_x^3\Delta_x^3] = 2(a\Delta)a_x^3\Delta_x^3 = 2t_x^6.
\]

Der Coefficient dieses Gliedes der Reihe (1) ist $\dfrac{\binom{4}{1}\binom{4}{1}}{\binom{8}{1}} = 2$; also erhalten wir:

\weq{I}{a_y^4\Delta_x^4 - a_x^4\Delta_y^4 = 4t_x^3 t_y^3(yx);}

daraus ergeben sich durch Polarisation nach $y$ die zwei weitern Formeln:

\weq{II}{a_y^3 a_x \Delta_x^4 - a_x^4 \Delta_y^3\Delta_x = 3t_x^4 t_y^2(yx)}

\weq{III}{a_y^2 a_x^2\Delta_x^4 - a_x^4\Delta_y^2\Delta_x^2 = 2t_x^5 t_y(yx),}
%CONT>>


% ===== pdf idx 198, printed 184.0 =====
% !PRINTED 184
\weberpage{184}

%<<CONT
wobei also:

\weq{IV}{(a\Delta)a_x^3\Delta_x^3 = t_x^6.}

Auf den rechten Seiten dieser vier Formeln befinden sich Polaren von $t$, von der nullten angefangen bis zur dritten. Je nachdem wir also für $y$ die Symbole $a$, $\Delta$, $t$ in dieselben eintragen, erhalten wir der Reihe nach die gewünschten Ueberschiebungen, ausgedrückt in Symbolen von $a$ und $\Delta$ allein.

164. \textit{Die ersten Ueberschiebungen.} Hiezu benutzen wir Formel III, wir setzen darin zunächst $y_1 = b_2$, $y_2 = -b_1$ und multipliciren mit $b_x^2$; dann erhält man:

\weq{1}{(ab)_x^2 a_x^2 b_x^2 \cdot \Delta - f\cdot(b\Delta)^2 b_x^2\Delta_x^2 = 2(t,f) = \Delta^2 - \frac{i}{6}\cdot f^2,}

sodann ersetzen wir $y_1$ durch $\Delta'$, $y_2$ durch $\Delta_2'$ und bekommen:

\weq{2}{(a\Delta')^2 a_x^2 \Delta_x'^2 \cdot \Delta - f\cdot(\Delta\Delta')^2\Delta_x^2\Delta_x'^2 = 2(t,\Delta) = \frac{1}{3} f(i\Delta - jf).}

Hiezu tritt als erste Ueberschiebung von $t$ über sich selbst:

\weq{3}{(t,t) = 0.}

165. \textit{Die zweiten Ueberschiebungen.} Wir ersetzen in Formel (II) zuerst $y$ durch $b$, dann $y$ durch $\Delta'$; beide Male erhalten wir als Resultat null; daher

\weq{4}{(t,f)^2 = 0}

\weq{5}{(t,\Delta)^2 = 0.}

Substituiren wir aber für $y$ das Symbol $t'$, so entsteht:
\[
(at)^3 a_x t_x^3 \cdot \Delta - f\cdot(\Delta t)^3 \Delta_x t_x^3 = 3(t,t')^2 t_x^4 t_x'^4,
\]
oder
\[
(f,t)^3 \cdot \Delta - (\Delta,t)^3 f = (t,t')^2.
\]
Führen wir für die sofort zu berechnenden Ueberschiebungen $(f,t)^3$ und $(\Delta,t)^3$ ihre aus (7) und (8) hervorgehenden Ausdrücke ein, so kommt:
\[
\frac{1}{4}\{jf - i\Delta\}\Delta - \frac{1}{4}\left\{\frac{i^2}{6} f - j\Delta\right\}f = 3(t,t')^2,
\]
oder

\weq{6}{-\frac{1}{12}\left\{i\Delta^2 - 2j\Delta f + \frac{i^2}{6} f^2\right\} = (t,t')^2.}

166. \textit{Die dritten Ueberschiebungen.} Sie gehen direct aus der Combinante $P$ selbst hervor [Gleichung (I)], indem man der Reihe nach $y$ durch $b$, $\Delta'$, $t'$ ersetzt. Man erhält:

\weq{7}{4(t,f)^3 = i\Delta - jf}

\weq{8}{4(t,\Delta)^3 = j\Delta - \frac{i^2}{6} f}

\weq{9}{4(t,t')^3 = 0.}


% ===== pdf idx 199, printed 185.0 =====
% !PRINTED 185
\weberpage{185}

167. \textit{Die vierten und fünften Ueberschiebungen.} Wir benutzen zu diesen Bildungen eine andere Ausgangsform.

Gemäss der Reihenentwicklung (IX), \S~7 ist:

\[
(a\Delta)\,\Delta_x^2\,a_y^3\,\Delta_y = \sum \frac{\binom{2}{k}\binom{3}{k}}{\binom{6-k+1}{k}}\,(f,\Delta)^k_{y^{3-k+1}}\,(xy)^k,
\]
oder
\[
(a\Delta)\,a_y^3\,\Delta_y\,\Delta_x^2 = t_{y^4} + \frac{i}{6}\cdot f_{y^3}(yx).
\]

Ersetzen wir hier $y$ durch $b$, so folgt:

\[
(a\Delta)(ab)^3(\Delta b)\Delta_x^2 = (t,f)^4 + \frac{i}{6}(f,f)^3.
\]

Nun ist aber $(f,f)^3=0$, und ebenso verschwindet die linke Seite, da sie durch Vertauschung von $a$ mit $b$ ihr Zeichen ändert, daher:

\weq{10}{(t,f)^4 = 0,}
und ebenso
\weq{11}{(t,\Delta)^4 = 0,}

wie aus (10) direct durch den Aronhold'schen Process (vergl. auch \S~16) hervorgeht. Beide Resultate sind auch a priori einleuchtend. Denn diese Ueberschiebungen würden quadratische Covarianten erzeugen, die aber unmöglich durch $f, \Delta, t$ sich rational ausdrücken lassen.

Die vierte Ueberschiebung von $t$ über sich selbst gewinnt man wiederum aus der Formel (II), Nr. 163, wenn man darin zuerst $x$ mit $y$ vertauscht, also aus:

\[
a_y^4\,\Delta_x^3\,\Delta_y - a_x^3\,a_y\,\Delta_y^4 = 3\,t_y^4\,t_x^2\,(yx).
\]

Man substituirt $t$ für $y$ und erhält:

\[
(at)^4(\Delta t)\,\Delta_x^3\,t_x - a_x^3(at)(\Delta t)^4 t_x = 3\,(t,t')^4\,t_x^2\,t_x'^{\,2}.
\]

Die beiden Glieder links verschwinden einzeln wegen der Symbolfactoren $(at)^4$, resp. $(\Delta t)^4$, vermöge welcher dieselben durch Ueberschiebung mit $(f,t)^4=0$, resp. $(\Delta t)^4=0$ entstanden gedacht werden können.

Somit habeu wir die wichtige Relation % [sic]

\weq{12}{(t,t)^4 = 0.}

Durch sie wird $t$ als eine ganz specielle Form ihres Grades charakterisirt. Wir werden in \S~19 solche Formen eingehender studiren. Hier möge nur so viel bemerkt sein, dass gerade wegen dieser Relation $(t,t)^4=0$
%CONT>>


% ===== pdf idx 200, printed 186.0 =====
% !PRINTED 186
\weberpage{186}

%<<CONT
nach den in \S~20 entwickelten Gesetzen das ganze Formensystem einer solchen Form $t$ sich auf die vier Formen reducirt

\[
t,\ (t,t)^2,\ ((t,t)^2,t),\ (t,t)^6.
\]

Insbesondere müssen daher alle Invarianten Potenzen von $(t,t)^6$ sein, welche, wie der nächste Absatz zeigt, den Werth $\frac{1}{4}\left(j^2-\frac{i^3}{6}\right)$ besitzt.

168. \textit{Die fünften und sechsten Ueberschiebungen.} Die fünften Ueberschiebungen sind durchwegs null, da ja $\Delta$ und $f$ nur vom vierten Grade sind und $(t,t)^5$ sein Zeichen ändert bei Vertauschung von $t$ mit $t'$.

Es bleibt sonach nur noch $(t,t)^6$ zu berechnen. Wir bilden von

\[
t_x^6 = (a\Delta)\,a_x^3\,\Delta_x^3
\]

die sechste Polare und ersetzen $y$ durch $t$:

\[
(t,t)^6 = (a\Delta)(at)^3(\Delta t)^3 = -(at)^3(a\Delta)(t\Delta)^3.
\]

Der Ausdruck rechts kann auch betrachtet werden als vierte Ueberschiebung von $\Delta$ über $(f,t)^3=(at)^3a_x t_x^3$; daher hat man:

\[
(t,t)^6 = -((f,t)^3,\Delta)^4 = -\frac{1}{4}\big((i\Delta-jf),\Delta\big)^4 = \frac{1}{4}\left(j^2-\frac{i^3}{6}\right).
\]

Wir werden alsbald sehen, dass die berechnete Grösse rechts nichts anderes ist als die Discriminante $R$ der Form $f=a_x^4$, abgesehen von einem Zahlenfactor.

\begin{center}
\S~16.

\gesperrt{Anwendung des Aronhold'schen Processes auf das Formensystem der Form vierten Grades.}
\end{center}

169. \textit{Wirkung des Processes auf die Grundformen des Systemes.} Wie die cubische Form eine Covariante $Q$ gleichen Grades besitzt, so haben wir auch im System der biquadratischen Form eine Covariante, die Hesse'sche Form $\Delta$, die gleichen Grades ist mit der Originalform $f$. Wir können daher in derselben Weise wie dort unter Einführung der Symbole $\Delta$, symbolische Producte, welche Covarianten von $f=a_x^4$ repräsentiren, dem Aronhold'schen Processe unterwerfen, und so wichtige Aufschlüsse über die Eigenschaften und den inneren Zusammenhang der Formen des Systems erhalten.

Ist $P$ irgend eine Covariante von $f=a_x^4$, so ist in unserm Falle der Aronhold'sche Process dargestellt durch:

\weq{I}{\delta P = \sum \frac{\partial P}{\partial \bar a_i}\,\bar\Delta_i,}
%CONT>>


% ===== pdf idx 201, printed 187.0 =====
% !PRINTED 187
\weberpage{187}

%<<CONT
wobei $\bar a_i$ und $\bar\Delta_i$ die unsymbolischen Co\"efficienten von $f$ resp. $\Delta$ bedeuten. Ersetzen wir nun in dieser Gleichung (I) die Covariante $P$ der Reihe nach durch $f, \Delta, i, t, j$, so erhalten wir:

\weq{1}{\delta f = \Delta;}

\weq{2}{\delta \Delta = \delta(f,f)^2 = 2(f,\delta f)^2 = 2(f,\Delta)^2 = \frac{i}{3}\cdot f.}

Wir haben also hier wiederum die wechselseitige Beziehung, die wir schon mehrmals bei analogen Untersuchungen betont haben. Ebenso wie $\delta f$ auf $\Delta$ führt, so gelangt man umgekehrt durch $\delta\Delta$ wieder auf $f$ zurück.

Es ist weiter:

\weq{3}{\delta i = \delta(f,f)^4 = 2(f,\delta f)^4 = 2(f,\Delta)^4 = 2j;}

\weq{4}{\delta t = \delta(f,\Delta) = (\delta f,\Delta)+(f,\delta\Delta) = (\Delta,\Delta)+i\cdot(f,f) = 0.}

Das letzte Resultat war a priori einzusehen, da $t$ als erste Ueberschiebung von $f$ und $\Delta$ eine Combinante dieser beiden Formen ist. (Vergl. \S~6.) Die allgemeinste Combinante des Systemes $f, \Delta$ ist die Form

\[
C = a_y^4\Delta_x^4 - a_x^4\Delta_y^4,
\]

wie wir bei Gelegenheit der Definition von Combinanten \S~6 erwähnt haben.

Da aber nach Gleichung Nr. 163 (I), $t$ die einzige Elementar-
%CONT>>


% ===== pdf idx 202, printed 188.0 =====
% !PRINTED 188
\weberpage{188}

%<<CONT
covariante von $C$ ist, so treten beide Combinanten $t$ wie $C$ gleichberechtigt im System auf. Endlich ist:

\weq{5}{\delta j = \delta(f,\Delta)^4 = (\delta f,\Delta)^4+(f,\delta\Delta)^4 = (\Delta,\Delta)^4+\frac{i}{3}(f,f)^4 = \frac{i^2}{6}+\frac{i^2}{3} = \frac{i^2}{2}.}

Damit ist an den Grundformen des Systemes die Wirkung des $\delta$-Processes studirt. Man kann die eben gewonnenen Resultate wiederum verwenden, aus einfacheren Ueberschiebungen schwierigere zu berechnen. So erhalten wir z.\,B. den Werth von $(\Delta,\Delta)^4$, wenn wir

\[
(f,\Delta)^2 = \frac{i}{6}f
\]

dem $\delta$-Process unterwerfen. Es ergiebt sich:

\[
(\Delta,\Delta)^2 + \frac{i}{3}(f,f)^2 = \frac{j}{3}f + \frac{i}{6}\Delta,
\]
also
\[
(\Delta,\Delta)^2 = \frac{1}{3}jf - \frac{1}{6}i\Delta,
\]

was mit dem früher gewonnenen Resultat, Nr. 158, (3), übereinstimmt.

170. \textit{Die Combinanten der Form $\varphi=kf+\lambda\Delta$.} Wir haben seinerzeit eine Form $C$ als Combinante zweier binärer Formen $\chi$ und $\psi$ definirt, sobald sie der Differentialgleichung

\[
\delta C = \sum \frac{\partial C}{\partial \bar\chi_i}\cdot \psi_i = 0
\]

genügt. Ist nun $P$ irgend eine Combinante von $f=a_x^4$ und $\Delta=\Delta_x^4$, also, wenn wir uns in $P$ die Co\"efficienten von $\Delta$ durch ihre Werthe in den Co\"efficienten von $f$ ersetzt denken, eine Combinante von $f$, so besitzt die Form

\[
\varphi = kf + \lambda\Delta
\]

eine entsprechende Combinante $P_\varphi$, welche man erhält, wenn man in $P$ die Co\"efficienten $\bar a_i$ von $f$ durch die Co\"efficienten $k\bar a_i+\lambda\bar\Delta_i$ von $\varphi$ ersetzt. Insbesondere entspricht der einzigen Fundamentalcombinante $t=(f,\Delta)$ eine Form $t_\varphi$ von $\varphi$, aus der alle Combinanten dieser Form $\varphi$ entspringen. Die Berechnung derselben ist also zurückgeführt auf die Berechnung von $t_\varphi$. Nach den allgemeinen Entwicklungen in \S~6 hat man aber:

\weq{II}{t_\varphi = g^\varrho\cdot t.}

Hiebei ist $\varrho$ eine noch zu bestimmende ganze Zahl; ferner ist

\[
g = kB - \lambda A,
\]

worin $A$ und $B$ die Co\"efficienten in der Gleichung

\[
\Delta_{kf+\lambda\Delta} = Af+B\Delta,
\]

sind. Es handelt sich also zunächst darum, $\Delta_{kf+\lambda\Delta}=\Delta_\varphi$ zu berechnen.

171. \textit{Berechnung von $\Delta_\varphi$.} Wir ermitteln diese Form direct durch Ueberschiebung. Man erhält:

\[
\begin{aligned}
\Delta_\varphi &= (kf+\lambda\Delta,\,kf+\lambda\Delta)^2\\
&= k^2\cdot(f,f)^2 + 2k\lambda\cdot(f,\Delta)^2 + \lambda^2\cdot(\Delta,\Delta)^2\\
&= k^2\cdot\Delta + \frac{i}{3}f\cdot k\lambda + \lambda^2\left(\frac{j}{3}f-\frac{i}{6}\Delta\right),
\end{aligned}
\]
also:
\[
\Delta_\varphi = \Delta\left\{k^2-\frac{i}{6}\lambda^2\right\} + f\left\{\frac{k\lambda i}{3}+\frac{j}{3}\lambda^2\right\}.
\]

Es ist sehr bemerkenswerth, dass die Coefficienten von $\Delta$ und $f$ auf der rechten Seite dieser Gleichung nichts anderes sind, als die
%CONT>>


% ===== pdf idx 203, printed 189.0 =====
% !PRINTED 189
\weberpage{189}
%<<CONT
ersten Differentialquotienten der Form
\weq{III}{g = k^3 - \frac{i}{2}\,k\lambda^2 - \frac{j}{3}\,\lambda^3}
nach $k$ beziehungsweise $\lambda$. Die Differentiation ergiebt nämlich unmittelbar
\[
\frac{1}{3}\,\frac{\partial g}{\partial k} = g_1 = k^2 - \frac{i}{6}\lambda^2
\]
\[
\frac{1}{3}\,\frac{\partial g}{\partial \lambda} = g_2 = -\frac{k\lambda i}{3} - \frac{j}{3}\lambda^2 .
\]
Die Form $g$ selbst aber ist eine cubische Resolvente der Gleichung
vierten Grades $f = a_x^4 = 0$, auf welche die Lagrange'sche Methode der
Auflösung dieser Gleichung direct führt*). Wir können also die Form
$\varDelta_\varphi$ einfacher darstellen durch
\[
\varDelta_\varphi = g_1\varDelta - g_2 f = B\varDelta + Af,
\]
und die Combinante $t_\varphi$ durch:
\[
t_\varphi = \left(k^3 - \frac{i}{2}\,k\lambda^2 - \frac{j}{3}\,\lambda^3\right)\cdot t .
\]
Denn die Constante $\varrho$ muss nothwendig den Werth $1$ besitzen,
da $t$ nur vom dritten Grade in den Coefficienten von $f$ und demnach
$t_\varphi$ auch nur vom dritten Grade in $k$ und $\lambda$ sein kann.

172. \textit{Berechnung der Formen $i_\varphi$ und $j_\varphi$.} Die letzten Betrachtungen
haben uns veranlasst, die beiden Formen $\varDelta$ und $t$ des Systemes von
$\varphi = kf + \lambda\varDelta$ zu berechnen. Wir wollen zur Vervollständigung auch
noch die beiden übrigen Formen $i$ und $j$ dieser Form $\varphi$ ermitteln,
und schlagen zu dem Zwecke wiederum den Weg der directen Ueberschiebung ein. Gemäss der Definition von $i$ als vierte Ueberschiebung
der Grundform über sich selbst wird:
\[
i_\varphi = (kf+\lambda\varDelta,\ kf+\lambda\varDelta)^4 = k^2(f,f)^4 + 2k\lambda(f,\varDelta)^4 + \lambda^2(\varDelta,\varDelta)^4
\]
\[
= k^2 i + 2k\lambda j + \lambda^2\,\frac{i^2}{6}.
\]
Genau dieselbe Form erhält man als quadratische (also Hesse'sche)
Covariante $\tau$ der cubischen Form $g$. Da nämlich
\[
g = k^3 - \frac{i}{2}\,k\lambda^2 - \frac{j}{3}\,\lambda^3,
\]

*) Sind $x_1\,x_2\,x_3\,x_4$ die Wurzeln von $f=0$, dann ist jene Gleichung, welche
die Grössen $y_1 = x_1x_2+x_3x_4$, $y_2 = x_1x_3+x_2x_4$, $y_3 = x_1x_4+x_2x_3$ zu Wurzeln
hat, cubische Resolvente von $f$. Aus ihr geht durch lineare Transformation die
Form $g$ hervor.
%CONT>>


% ===== pdf idx 204, printed 190.0 =====
% !PRINTED 190
\weberpage{190}
%<<CONT
so wird die Determinante der zweiten Differentialquotienten:
\[
\frac{1}{2}\,\tau_g = \frac{1}{36}
\begin{vmatrix}
\dfrac{\partial^2 g}{\partial k^2} & \dfrac{\partial^2 g}{\partial k\,\partial \lambda} \\
\dfrac{\partial^2 g}{\partial k\,\partial \lambda} & \dfrac{\partial^2 g}{\partial \lambda^2}
\end{vmatrix}
=
\begin{vmatrix}
k, & \dfrac{-i}{6}\lambda \\
\dfrac{-i}{6}\lambda, & -\left(\dfrac{ik}{6}+\dfrac{j}{3}\lambda\right)
\end{vmatrix}
\]
\[
\tau_g = -\frac{1}{3}\left\{ i k^2 + 2k\lambda j + \frac{i^2}{6}\lambda^2 \right\};
\]
daher
\[
i_\varphi = -3\tau_g .
\]
Wir erhalten endlich:
\[
j_\varphi = (kf+\lambda\varDelta,\ g_1\varDelta - g_2 f)^4 = k g_1 j - k g_2 i + \lambda g_1\,\frac{i^2}{6} - \lambda g_2 j
\]
\[
= g_1\left(kj + \frac{\lambda i^2}{6}\right) - g_2(ki+\lambda j).
\]
Auch dieser Ausdruck rechts lässt sich als cubische Covariante $q$ der
Form $g$ herstellen; denn es ist
\[
q_g = (g,\ \tau_g) = \frac{1}{6}
\begin{vmatrix}
3g_1, & 3g_2 \\
-\dfrac{2}{3}(ik+\lambda j), & -\dfrac{2}{3}\left(kj+\dfrac{i\lambda}{6}\right)
\end{vmatrix}
\]
\[
= -\frac{1}{3}\left\{ g_1\left(kj+\frac{i\lambda}{6}\right) - g_2(ik+\lambda j) \right\}.
\]
Daher:
\[
j_\varphi = -3(g,\ \tau_g) = -3 q_g .
\]
Setzt man in die Relation, welche zwischen den In- und Covarianten
 der cubischen Form existirt [vgl. Nr. 148 (I)]
\[
q_g^2 = -\frac{1}{2}(\tau_g^3 + R_g g^2)
\]
die Werthe $q_g$, $\tau_g$, $R_g$, $g$ ein, ausgedrückt in den Formen $j$ und $i$,
so erhält man nach leichter Umformung auch noch
\[
i_\varphi^3 - 6 j_\varphi^2 = g^2(i^3 - 6j^2)
\]
eine Relation, welche nach Nr. 170 (II) diese Verbindung als Combinante
 charakterisirt. Das ganze Formensystem der Form $\varphi$ ist also
dargestellt durch:
\[
\varphi = kf+\lambda\varDelta, \quad \varDelta_\varphi = g_1\varDelta - g_2 f, \quad t_\varphi = g\cdot t,
\]
\[
i_\varphi = -3\tau_g, \quad j_\varphi = -3q_g,
\]
wobei $g = k^3 - \dfrac{ik\lambda^2}{2} - \dfrac{j}{3}\lambda^3$, und $g_1 = \dfrac{1}{3}\dfrac{\partial g}{\partial k}$, $g_2 = \dfrac{1}{3}\dfrac{\partial g}{\partial \lambda}$ ist.


% ===== pdf idx 205, printed 191.0 =====
% !PRINTED 191
\weberpage{191}
\begin{center}
\S~17.\\[2pt]
\textbf{Die Gleichung vierten Grades.}
\end{center}

173. \textit{Die Discriminante von $f=a_x^4$.} Ehe wir uns mit der Auflösung
 der Gleichung vierten Grades und mit den Eigenschaften ihrer
Wurzeln eingehender beschäftigen, wollen wir zunächst jene Invariante
der Form vierten Grades berechnen, welche als Discriminante bezeichnet
wird. Sie muss sich nach dem Vorausgegangenen rational durch $i$
und $j$ darstellen lassen, da diese beiden Formen die einzigen selbstständigen
 Invarianten des Systemes sind.

Gesetzt nun, die Form $f=a_x^4$ besitze zwei gleiche Wurzeln und
sei dementsprechend dargestellt durch
\weq{1}{f = a_x^4 = \alpha_x^2\cdot p_x^2 ,}
wo $\alpha_x$ der lineare Doppelfactor, und $p_x^2$ irgend eine quadratische Form
ist, dann muss für diese Form die Discriminante identisch verschwinden
. Es muss also in diesem Falle zwischen den beiden Fundamentalinvarianten
 $i$ und $j$ eine Relation bestehen, die sich von selbst
einstellen muss, sobald wir $i$ und $j$ für $f=\alpha_x^2\cdot p_x^2$ berechnen. Beide
Formen werden in diesem Falle durch einen gewissen Parameter $\varrho$
sich darstellen lassen, durch dessen Elimination sich die gewünschte
Relation ergiebt.

Nun hat die zweite Ueberschiebung von $f=\alpha_x^2\cdot p_x^2$ über $\alpha_x^2$ den
Werth:
\weq{2}{(a\alpha)^2 a_x^2 = \frac{1}{6}(p\alpha)^2\cdot\alpha_x^2 = \varrho\cdot\alpha_x^2}
und daher erhalten wir, indem wir diese Gleichung einmal über
$f=b_x^4$ schieben, als erste Ueberschiebung von $f$ über $\alpha$:
\weq{3}{(ab)(a\alpha)^2 a_x b_x^3 = -\varrho\cdot\alpha\cdot(f,\ \alpha).}
Die linke Seite enthält den Reducenten $(ab)$ und muss daher auf eine
Ueberschiebung von $\varDelta$ über $\alpha$ führen. In der That erhält man durch
Anwendung des Identitätssatzes:
\[
(ab)(a\alpha)^2 a_x b_x^3 = (ab)(a\alpha)\,a_x b_x^2\{(b\alpha)a_x + (ab)\alpha_x\},
\]
oder, weil hier das erste Glied rechts durch Vertauschung von $a$ mit
$b$ nur sein Zeichen ändert und sonach verschwindet:
\weq{4}{(ab)(a\alpha)^2 a_x b_x^3 = (ab)^2(a\alpha)\,a_x b_x^2\alpha_x = \alpha\cdot(\varDelta,\ \alpha).}
Subtrahirt man nun beide Gleichungen (3) und (4), so erhält man
nach Absonderung des Factors $\alpha_x$:
\weq{5}{(\varDelta+\varrho\cdot f,\ \alpha) = 0.}
Da nun die Invarianten $i$ und $j$ durch viermalige Ueberschiebung von
$\varDelta$ und $f$ über sich selbst und über einander entstehen, so brauchen
%CONT>>


% ===== pdf idx 206, printed 192.0 =====
% !PRINTED 192
\weberpage{192}
%<<CONT
wir nur, um die gewünschten Relationen für $i$ und $j$ zu erhalten,
die Gleichung (5) dreimal über $\varDelta$ und über $f$, oder besser über die Combination
 $\varDelta+\mu f$ zu schieben, wo $\mu$ jeden beliebigen Werth hat. Wir
erhalten dann zunächst:
\[
U = ((\varDelta+\varrho f,\ \alpha),\ \varDelta+\mu f)^3 = ((P,\ \alpha),\ Q)^3 = 0.
\]
Ersetzen wir hier einen Moment $\alpha$ durch $y$ und entwickeln das so
sich ergebende Polarenglied: $(PQ)^3 P_y Q_x$ in eine Reihe, so kommt,
wenn wir darin wiederum für $y$ das Symbol $\alpha$ einführen:
\[
U = 0 = ((\varDelta+\varrho f,\ \varDelta+\mu f)^3,\ \alpha) + \frac{1}{2}(\varDelta+\varrho f,\ \varDelta+\mu f)^4\cdot\alpha,
\]
oder, weil alle dritten Ueberschiebungen $(\varDelta,\varDelta)^3$, $(\varDelta,f)^3$, $(f,f)^3$ verschwinden
:
\[
(\varDelta+\varrho f,\ \varDelta+\mu f)^4 = 0,
\]
d.\,h.
\weq{6}{\frac{i^2}{6} + (\varrho+\mu)j + \varrho\mu\cdot i = 0.}
Da diese Identität für jeden Werth $\mu$ gilt, so müssen die Beziehungen
bestehen
\[
\frac{i^2}{6} + \varrho\cdot j = 0
\]
\[
i\varrho + j = 0.
\]
Durch Elimination des Parameters $\varrho$ erhält man
\[
R_f = j^2 - \frac{i^3}{6} = 0
\]
und in dieser Relation muss die linke Seite als Discriminante der
Form $f$ angesehen werden, da diese Relation nur entstanden ist unter
der Bedingung des Vorhandenseins einer Doppelwurzel; sie ist in der
That auch vom Grade $2(n-1)=6$ in den Coefficienten von $f$.

\textit{Anmerkung.} Berechnet man die Discriminante $R_g$ der cubischen
Resolvente
\[
g = k^3 - \frac{i}{2}\,k\lambda^2 - \frac{j}{3}\,\lambda^3
\]
nach den Regeln der niederen Algebra (vgl. auch Bd. I Nr. 183), so
ergiebt sich
\[
R_g = \frac{1}{9}\left(j^2 - \frac{i^3}{6}\right).
\]
Ebenso haben wir gesehen, dass die Invariante
\[
(t,\ t)^6 = \frac{1}{4}\left(j^2 - \frac{i^3}{6}\right).
\]

174. \textit{Die Cayley'sche Auflösung der Gleichung vierten Grades.} Die
invariantentheoretische Auflösung der Gleichung
\[
f = \bar a_0 x_1^4 + 4\bar a_1 x_1^3 x_2 + 6\bar a_2 x_1^2 x_2^2 + 4\bar a_3 x_1 x_2^3 + \bar a_4 x_2^4 = 0
\]
%CONT>>


% ===== pdf idx 207, printed 193.0 =====
% !PRINTED 193
\weberpage{193}
%<<CONT
knüpft Cayley an die in Nr. 162 aufgestellte Relation
\[
2t^2 = \frac{i}{2}\, f^2\, \Delta - \Delta^3 - \frac{j}{3}\, f^3,
\]
die wir auch in der Form schreiben können
\[
-2t^2 = \Delta^3 - \frac{i}{2}\, f^2\, \Delta - \frac{j}{3}\, (-f)^3.
\]
Vergleicht man nun die rechte Seite dieser Relation mit der cubischen Resolvente
\[
g = k^3 - \frac{i}{2}\, \lambda^2 k - \frac{j}{3}\, \lambda^3
\]
der Gleichung $f = 0$, so erkennt man, dass dieselbe daraus hervorgeht, wenn man $k$ durch $\Delta$ und $\lambda$ durch $-f$ ersetzt. Man hat also:
\[
-2t^2 = g(\Delta,\, -f)
\]
oder, wenn $k_1$, $k_2$, $k_3$ die drei Wurzeln von $g=0$ sind,
\weq{1}{-2t^2 = (\Delta+k_1f)(\Delta+k_2f)(\Delta+k_3f).}
Ist nun $f$ eine allgemeine Form ihres Grades, sind also ihre vier Wurzeln von einander verschieden, so ist die Discriminante $R$ von $f$ nicht null; folglich sind in diesem Falle keine zwei der drei Wurzeln $k_1$, $k_2$, $k_3$ einander gleich (da $R$ auch Discriminante der Resolvente $g$ ist); anderntheils können $f$ und $\Delta$ im Allgemeinen keinen Factor gemeinsam haben, wenn $R \gtrless 0$*)\footnote{Vgl. auch Nr.~177, Anmerkung.}.

Die drei Factoren der Gleichung (I) sind also in Bezug auf ihre linearen Factoren von einander verschieden, und da die linke Seite ein völliges Quadrat ist, so muss jeder der drei Factoren rechts gleichfalls das Quadrat einer Form zweiten Grades sein.

Wir können daher setzen
\weq{2}{\left.\begin{aligned}
\Delta + k_1 f &= \varphi^2 \\
\Delta + k_2 f &= \psi^2 \\
\Delta + k_3 f &= \chi^2
\end{aligned}\right\}}
wo $\varphi$, $\psi$ und $\chi$ drei quadratische Formen sind.

Alsdann wird die Covariante
\[
t = \frac{1}{\sqrt{-2}} \cdot \varphi \cdot \psi \cdot \chi.
\]
Die drei quadratischen Formen $\varphi$, $\psi$, $\chi$ sind aber conjugirte Formen (vgl. Nr. 137); denn die Functionaldeterminanten je zweier derselben sind immer der dritten proportional. So ist z. B.
%CONT>>


% ===== pdf idx 208, printed 194.0 =====
% !PRINTED 194
\weberpage{194}
%<<CONT
\[
\begin{aligned}
(\Delta+k_2f,\ \Delta+k_3f) &= (\Delta,\Delta)+(k_2-k_3)(f,\Delta)+k_2k_3(f,f) \\
&= (k_2-k_3)\cdot t \\
&= \frac{k_2-k_3}{\sqrt{-2}} \cdot \varphi \cdot \psi \cdot \chi.
\end{aligned}
\]
Anderntheils ist
\[
(\psi^2,\ \chi^2) = (\psi,\chi)\cdot\psi\cdot\chi.
\]
Durch Comparation beider Gleichungen ergiebt sich
\weq{3}{\left.\begin{aligned}
(\psi,\chi) &= \frac{k_2-k_3}{\sqrt{-2}}\,\varphi \\
\text{und analog}\quad (\chi,\varphi) &= \frac{k_3-k_1}{\sqrt{-2}}\,\psi \\
(\varphi,\psi) &= \frac{k_1-k_2}{\sqrt{-2}}\,\chi
\end{aligned}\right\}.}
Die Covariante $t$ ist also eine ganz specielle Form sechsten Grades, welche unter Adjunction der drei Wurzeln einer cubischen Gleichung in drei conjugirte quadratische Factoren gespalten werden kann.

175. Gerade dieser Umstand gestattet aber die Wurzeln der Gleichung $f=0$ in einfacher Weise zu berechnen. Durch Subtraction je zweier der Gleichungen (2) erhalten wir nämlich
\weq{4}{\left.\begin{aligned}
(k_1-k_2)f &= \varphi^2-\psi^2 = (\varphi-\psi)(\varphi+\psi) \\
(k_2-k_3)f &= \psi^2-\chi^2 = (\psi-\chi)(\psi+\chi) \\
(k_3-k_1)f &= \chi^2-\varphi^2 = (\chi-\varphi)(\chi+\varphi)
\end{aligned}\right\}.}
Wir schliessen aus diesen Gleichungen (4), dass je eine Summe und Differenz je zweier solcher quadratischer Formen alle Wurzeln von $f$ enthalten, dass daher die Differenzen
\[
\begin{aligned}
&\varphi-\psi \quad \text{und} \quad \psi-\chi \\
\text{oder}\quad &\psi-\chi \quad \text{und} \quad \chi-\varphi \\
\text{oder}\quad &\chi-\varphi \quad \text{und} \quad \varphi-\psi
\end{aligned}
\]
oder die entsprechenden Summenpaare stets einen gemeinsamen linearen Factor von $f$ besitzen. Denn enthielten sie keinen solchen gemeinsamen Factor, so müsste sich einmal $f$ darstellen lassen z. B. durch
\[
(k_1-k_2)f = (\varphi-\psi)(\varphi+\psi),
\]
dann aber auch durch
\[
\varrho \cdot f = (\varphi-\psi)(\varphi-\chi).
\]
Also müsste bis auf eine Constante $\chi = -\psi$ sein, und analog $\varphi = -\chi$, d. h. alle drei quadratischen Formen müssten übereinstimmen und demnach $t$ Doppelwurzeln besitzen. Dann verschwindet aber die Discriminante von $t$, welche nur eine Potenz der Discriminante $R$ von $f$ sein kann,
%CONT>>


% ===== pdf idx 209, printed 195.0 =====
% !PRINTED 195
\weberpage{195}
%<<CONT
wie wir in Nr. 167 gesehen haben. Nach unserer Voraussetzung ist aber $R$ von null verschieden, und folglich haben in der That je zwei solche oben angeführte Summen- oder Differenzenpaare einen gemeinsamen Factor.

Nun hatten wir früher (Nr. 128) den Satz erhalten: Besitzen zwei quadratische Formen einen gemeinsamen linearen Factor, so enthält ihre Functionaldeterminante denselben quadratisch. Also ist in der Relation
\[
\sqrt{-2}\cdot(\varphi-\psi,\ \varphi-\chi) = (k_3-k_1)\psi+(k_1-k_2)\chi+(k_2-k_3)\varphi
\]
die rechte Seite das vollständige Quadrat des gesuchten linearen Factors von $f$. Wir können diese rechte Seite auch in die Determinantenform bringen
\[
-D = \begin{vmatrix} 1 & k_1 & \varphi \\ 1 & k_2 & \psi \\ 1 & k_3 & \chi \end{vmatrix}.
\]
Um den Factor selbst zu erhalten, brauchen wir nur die erste Polare zu bilden; dann stellt
\[
-D_y = \begin{vmatrix} 1 & k_1 & \varphi_y \\ 1 & k_2 & \psi_y \\ 1 & k_3 & \chi_y \end{vmatrix}
\]
für $y_2=1$ und $y=0$ einen linearen Factor von $f$ dar. Nun hat bereits Euler gezeigt, dass, wenn $x_1=\sqrt{t}+\sqrt{u}+\sqrt{v}$ eine Lösung der Gleichung vierten Grades ist, die drei andern Lösungen daraus dadurch erhalten werden, dass man je zwei der drei Grössen $\sqrt{t}$, $\sqrt{u}$, $\sqrt{v}$ mit negativem Vorzeichen nimmt. Die vier linearen Factoren der Gleichung vierten Grades sind also:
\[
\begin{vmatrix} 1, & k_1, & (+\sqrt{\Delta+k_1f})_y \\ 1, & k_2, & (+\sqrt{\Delta+k_2f})_y \\ 1, & k_3, & (+\sqrt{\Delta+k_3f})_y \end{vmatrix},
\quad
\begin{vmatrix} 1, & k_1, & (+\sqrt{\Delta+k_1f})_y \\ 1, & k_2, & (-\sqrt{\Delta+k_2f})_y \\ 1, & k_3, & (-\sqrt{\Delta+k_3f})_y \end{vmatrix},
\]
\[
\begin{vmatrix} 1, & k_1, & (-\sqrt{\Delta+k_1f})_y \\ 1, & k_2, & (+\sqrt{\Delta+k_2f})_y \\ 1, & k_3, & (-\sqrt{\Delta+k_3f})_y \end{vmatrix},
\quad
\begin{vmatrix} 1, & k_1, & (-\sqrt{\Delta+k_1f})_y \\ 1, & k_2, & (-\sqrt{\Delta+k_2f})_y \\ 1, & k_3, & (+\sqrt{\Delta+k_3f})_y \end{vmatrix}.
\]
\textit{Anmerkung.} Eine zweite und dritte Methode der Auflösung einer Gleichung vierten Grades wird in den nächsten Paragraphen gegeben werden.

176. \textit{Bedingung, dass $f=0$ drei gleiche Wurzeln enthält.} An die Auflösung der Gleichung vierten Grades knüpfen sich naturgemäss die
%CONT>>


% ===== pdf idx 210, printed 196.0 =====
% !PRINTED 196
\weberpage{196}
%<<CONT
Fragen: Unter welchen Bedingungen wird es eintreten, dass zwei, drei, vier der eben ermittelten Wurzelgrössen übereinstimmen? Was den ersten und letzten Fall betrifft, so sind die Bedingungen hiefür bereits ermittelt.

Die Gleichung $f=0$ hat zwei gleiche Wurzeln, sobald die Discriminante
\[
j^2 - \frac{i^3}{6} = 0.
\]
Die Gleichung $f=0$ hat vier gleiche Wurzeln, wenn die Hesse'sche Form identisch verschwindet, also wenn
\[
\Delta = (ab)^2 a_x^2 b_x^2 = 0 \quad \text{(vgl. Nr. 55)}.
\]
Es bleibt also noch zu untersuchen, unter welcher Bedingung $f=0$ drei gleiche Wurzeln besitzt. Zu dem Zwecke wählen wir $f=a_x^4$ in einer Normalform, in welcher von vorn herein drei Wurzeln als gleich angenommen sind, nämlich
\[
f = 4x_1^3 x_2.
\]
Die Coefficienten $a_0\, a_1\, a_2\, a_3\, a_4$ sind also der Reihe nach: $0, 1, 0, 0, 0$. Trägt man die Werthe dieser Coefficienten in die Invarianten $i$ und $j$ ein, so findet sich:
\[
i = 0,\quad j = 0.
\]
Wir fragen uns nun, ob die beiden Bedingungen $i=0$, $j=0$ auch hinreichend sind, damit $f=0$ drei gleiche Wurzeln besitze. Da in Folge des Verschwindens dieser Invarianten auch die Discriminante $R=0$ wird, so hat $f$ sicher zwei gleiche Wurzeln und kann also in diesem Falle in der Form dargestellt werden
\[
f = a_0x_1^4 + 4a_1x_1^3x_2 + 6a_2x_1^2x_2^2.
\]
Bildet man aber für diese Form die Invariante $i$, so reducirt sie sich auf $a_2^2=0$, d. h. auch $a_2$ muss nach Voraussetzung verschwinden; dann ist aber $f$ von der Form
\[
f = a_0x_1^4 + 4a_1x_1^3x_2,
\]
d. h. $f=0$ besitzt drei gleiche Wurzeln, und $j=0$, $i=0$ sind also die nothwendigen und hinreichenden Bedingungen dafür.

177. \textit{Bedingung, dass $f=a_x^4$ ein vollständiges Quadrat ist.} Es kann nun endlich noch der Fall eintreten, dass $f$ zwei Paare gleicher Wurzeln enthält, oder mit anderen Worten in das Quadrat einer Form zweiten Grades zerfällt. Um auch hiefür die Kriterien zu gewinnen, nehmen wir $f$ in einer entsprechenden Normalform
\[
f = cx_1^2 x_2^2
\]
%CONT>>


% ===== pdf idx 211, printed 197.0 =====
% !PRINTED 197
\weberpage{197}
und bilden die Hesse'sche Covariante derselben. Sie wird (vgl. Nr. 54)
\[
\Delta=\frac{1}{18}\begin{vmatrix}cx_2{}^2 & 2cx_1x_2\\ 2cx_1x_2 & cx_1{}^2\end{vmatrix}=-\frac{1}{6}c^2x_1{}^2x_2{}^2=\varrho\cdot f.
\]
In diesem Falle ist sonach $\Delta$ proportional der Form $f$. Man kann auch sagen: Wenn $f$ ein vollständiges Quadrat ist, dann bestehen zwischen den Coefficienten von $f$ und $\Delta$ die Relationen:
\[
\overline{a_0}=\varrho\,\overline{\Delta_0},\quad \overline{a_1}=\varrho\,\overline{\Delta_1},\quad \overline{a_2}=\varrho\,\overline{\Delta_2},\quad \overline{a_3}=\varrho\,\overline{\Delta_3},\quad \overline{a_4}=\varrho\,\overline{\Delta_4}.
\]
Wir fragen uns nun: Ist umgekehrt, wenn $f=\varrho\cdot\Delta$, diese Bedingung hinreichend, damit $f$ zwei Paare gleicher Wurzeln enthält? Um die Frage zu entscheiden, bilden wir
\[
(yx)^2\cdot\Delta_x^4=(yx)^2(ab)^2a_x^2b_x^2=(a_xb_y-b_xa_y)^2a_x^2b_x^2=2\begin{vmatrix}f, & f_y\\ f_y, & f_{y^2}\end{vmatrix}.
\]
Ist nun $x=\alpha$ irgend eine Wurzel von $f$, so verschwindet in der Gleichung:
\weq{I}{(yx)^2\Delta_x^4=2\begin{vmatrix}f, & f_y\\ f_y, & f_{y^2}\end{vmatrix},}
wenn wir dieselbe für $x$ eintragen, zunächst $f$ selbst, dann aber auch $\Delta$ wegen der Proportionalität mit $f$. Sie reducirt sich also auf:
\[
0=f_y^2\quad\text{oder}\quad f_y=0,
\]
d.\,h. auch die erste Polare von $f$ verschwindet, wenn $x=\alpha$, für jeden Werth von $y$. Dies ist aber nur möglich, wenn in derselben
\[
f_y=y_1\frac{\partial f}{\partial x_1}+y_2\frac{\partial f}{\partial x_2}
\]
die Differentialquotienten für $x=\alpha$ einzeln verschwinden. Dann aber ist diese Wurzel $x=\alpha$ eine Doppelwurzel. Wenn aber jede Wurzel $\alpha$ von $f$ Doppelwurzel ist, so ist $f$ ein vollständiges Quadrat.

\textit{Anmerkung.} Aus der eben benutzten Relation (I) geht direct hervor, dass $\Delta$ und $f$ keine Wurzeln gemein haben können, wenn $R\gtrless 0$. Denn ersetzen wir darin $x$ der Reihe nach durch die vier Wurzeln $\alpha_1\,\alpha_2\,\alpha_3\,'\alpha_4$ von $f$ und multipliciren die vier Identitäten, so kommt: % [sic] Komma vor $\alpha_4$ wie gedruckt
\[
\Delta(\alpha_1)\cdot\Delta(\alpha_2)\cdot\Delta(\alpha_3)\cdot\Delta(\alpha_4)=16\cdot\frac{f_y^2(\alpha_1)\cdot f_y^2(\alpha_2)\cdot f_y^2(\alpha_3)\cdot f_y^2(\alpha_4)}{(y\alpha_1)^2\cdot(y\alpha_2)^2\cdot(y\alpha_3)^2\cdot(y\alpha_4)^2}
\]
oder für $y_2=0$, $y_1=1$ erhalten wir (vgl. auch Bd. I Nr. 168, 173)
\[
(-1)^{4\cdot4}R_{\Delta,f}=16\cdot\varrho\left\{\left(\frac{\partial f}{\partial x}\right)_{x=\alpha_1}\cdot\left(\frac{\partial f}{\partial x}\right)_{x=\alpha_2}\cdot\left(\frac{\partial f}{\partial x}\right)_{x=\alpha_3}\cdot\left(\frac{\partial f}{\partial x}\right)_{x=\alpha_4}\right\}^2
\]
oder
\[
R_{\Delta,f}=16\cdot\varrho R_{f_1,f}^2=\mu\cdot R^2,\quad\text{d.\,h.}
\]
\glqq die Resultante von $\Delta$ und $f$ ist bis auf einen Zahlenfactor gleich dem Quadrat der Discriminante $R$ von $f$.\grqq{}


% ===== pdf idx 212, printed 198.0 =====
% !PRINTED 198
\weberpage{198}
\S~18. \textbf{Normalform von $f=a_x^4$; Transformation in die Normalform.}

178. \textit{Das Formensystem der Normalform.} Bei den Formen vierter Ordnung ist vorzugsweise ein Ausdruck derselben als canonische Form oder Normalform bezeichnet worden, nämlich jener, in welchem alle ungeraden Potenzen von $x$ fehlen, also
\[
f=x_1{}^4+6mx_1{}^2x_2{}^2+x_2{}^4.
\]
Wir studiren zunächst die algebraischen Ausdrücke des Systemes dieser Form, um alsdann umgekehrt aus dessen Eigenschaften die Mittel zu gewinnen, eine beliebige Form vierten Grades auf ihre Normalform zu transformiren.

Da die Coefficienten von $f$ der Reihe nach die Werthe haben
\[
1,\ 0,\ m,\ 0,\ 1,
\]
so wird die Hesse'sche Covariante derselben repräsentirt durch
\[
\Delta=2\begin{vmatrix}x_1{}^2+mx_2{}^2, & 2mx_1x_2\\ 2mx_1x_2 & mx_1{}^2+x_2{}^2\end{vmatrix}=2m(x_1{}^4+x_2{}^4)+2(1-3m^2)x_1{}^2x_2{}^2.
\]
Die rechte Seite dieser Gleichung enthält gleichfalls nur gerade Potenzen, also tritt auch $\Delta$ in der Normalform auf. Die Invarianten
\[
i=2(\overline{a_0}\,\overline{a_4}-4\overline{a_1}\,\overline{a_3}+3\overline{a_2}{}^2),\quad j=6\begin{vmatrix}\overline{a_0} & \overline{a_1} & \overline{a_2}\\ \overline{a_1} & \overline{a_2} & \overline{a_3}\\ \overline{a_2} & \overline{a_3} & \overline{a_4}\end{vmatrix}
\]
erhalten für die canonische Form von $f$ die Werthe:
\[
i=2(1+3m^2),\quad j=6m(1-m)^2.
\]
Endlich die Functionaldeterminante $t=(f,\ \Delta)$ ergiebt sich in der Form
\[
t=\frac{1}{16}\begin{vmatrix}\dfrac{\partial f}{\partial x_1}, & \dfrac{\partial f}{\partial x_2}\\ \dfrac{\partial\Delta}{\partial x_1}, & \dfrac{\partial\Delta}{\partial x_2}\end{vmatrix}=2\begin{vmatrix}x_1{}^3+3mx_1x_2{}^2, & x_2{}^3+3mx_1{}^2x_2\\ mx_1{}^3+\frac{1}{2}(1-3m^2)x_1x_2{}^2, & mx_2{}^3+\frac{1}{2}(1-3m^2)x_1{}^2x_2\end{vmatrix}
\]
\[
=(1-9m^2)x_1x_2(x_1{}^4-x_2{}^4)=(1-9m^2)\cdot(x_1{}^2-x_2{}^2)(x_1{}^2+x_2{}^2)x_1x_2
\]
\[
=(1-9m^2)\varphi\cdot\chi\cdot\psi.
\]
Die Covariante $t$ zerfällt also unter Zugrundelegung der obigen Normalform von $f$ nicht nur a priori in ihre drei conjugirten quadratischen Formen, sondern diese quadratischen Formen stellen sich hiebei selbst in der einfachsten Normalform ein; insbesondere reducirt sich $\varphi$ auf das Product der Linearfactoren $x_1\cdot x_2$. Dieser Umstand legt die Frage nahe, ob nicht umgekehrt durch die Transformation
%CONT>>


% ===== pdf idx 213, printed 199.0 =====
% !PRINTED 199
\weberpage{199}
%<<CONT
\[
\begin{aligned}
y_1&=r_x\\
y_2&=s_x,
\end{aligned}
\]
wobei $r_x$ und $s_x$ die Linearfactoren der Form $\varphi$ einer in allgemeiner Darstellung gegebenen Form $f$ sind, diese Form $f$ auf die oben erwähnte Normalform gebracht werden kann.

179. \textit{Transformation der Form $f$ in die Normalform.} In der That erreichen wir durch diese Transformation das gewünschte Ziel. Zunächst geht aus der Annahme $\varphi=r_x\cdot s_x$ auch die Gestalt der beiden andern conjugirten Formen hervor, nämlich (vgl. Nr. 137)
\weq{1}{\psi=r_x^2-s_x^2,\ \chi=r_x^2+s_x^2}
und wir wissen aus der Auflösung der Gleichung vierten Grades, dass, wenn wir setzen:
\weq{2}{\varphi^2=\Delta+k_1f,}
die Quadrate von $\psi$ und $\chi$ dargestellt sind durch
\weq{3}{\psi^2=\Delta+k_2f}
\weq{4}{\chi^2=\Delta+k_3f,}
wo $k_1$, $k_2$, $k_3$ die Wurzeln der cubischen Resolvente $g=0$ sind. Nun können wir mit Hilfe der Identität
\[
a_x(rs)=(as)r_x-(ar)s_x
\]
den Ausdruck $f\cdot(rs)^4=a_x^4(rs)^4=[a_x(rs)]^4$ direct nach dem binomischen Lehrsatze entwickeln und erhalten
\[
f\cdot(rs)^4=(as)^4r_x^4-4(as)^3(ar)r_x^3s_x+6(as)^2(ar)^2r_x^2s_x^2-4(as)(ar)^3r_xs_x^3
\]
\[
+(ar)^4s_x^4.
\]
Wir haben also nur zu zeigen, dass in Folge der Beziehungen (1) (2) (3) und (4) das zweite und vierte Glied dieses Ausdruckes rechts verschwindet; dann ist in der That $f$ auf die gewünschte Normalform gebracht.

Zu dem Zwecke untersuchen wir die Ueberschiebungen $(\varphi,\ r_x^2)^2$ und $(\psi,\ r_xs_x)^2$. Beide sind null; denn nach Nr. 38 erhalten wir:
\[
(\varphi,\ r_x^2)^2=(r_x\cdot s_x,\ r_x\cdot r_x)^2=(rr)(sr)=0
\]
und
\[
(\psi,\ \varphi)^2=(r_x^2-s_x^2,\ r_x\cdot s_x)^2=(rr)(rs)-(sr)(ss)=0.
\]
Daraus folgt aber unmittelbar, dass auch:
\weq{5}{(\varphi^2,\ r_x^3\cdot s_x)^4=((\varphi^2,\ r_x^2)^2,\ r_xs_x)^2=0}
und
\weq{6}{(\psi^2,\ r_x^3\cdot s_x)^4=((\psi^2,\ \varphi)^2,\ r_x^2)^2=0}
ist, oder, indem wir beide Relationen combiniren, dass auch
\weq{7}{(\varphi^2-\psi^2,\ r_x^3s_x)^4=0.}


% ===== pdf idx 214, printed 200.0 =====
% !PRINTED 200
\weberpage{200}
Nun ist aber wegen der Gleichungen (2) und (3)
\[
\varphi^2-\psi^2=(k_1-k_2)f;
\]
also ist auch gemäss der Relation (7):
\weq{8}{(k_1-k_2)(f,\ r^3s)^4=(k_1-k_2)(ar)^3(as)=0}
und ebenso folgt, wenn wir in (8) $r$ mit $s$ vertauschen,
\[
(f,\ s^3r)^4=(ar)(as)^3=0.
\]
Die Coefficienten des zweiten und vierten Gliedes in der binomischen Entwicklung von $f\cdot(rs)^4$ verschwinden also, und daher hat man in der That:
\weq{A}{f\cdot(rs)^4=(ar)^4s_x^4+6(ar)^2(as)^2r_x^2s_x^2+(as)^4r_x^4.}
\textit{Anmerkung über die Verwerthung der Normalform bei Auflösung der Gleichung vierten Grades.} Die eben zum Abschluss gebrachten Untersuchungen geben eine weitere Methode, die Gleichung vierten Grades aufzulösen. Man berechnet zunächst die cubische Resolvente derselben
\[
g=k^3-\frac{i}{2}k-\frac{j}{3}
\]
und bestimmt die drei Wurzeln $k_1\,k_2\,k_3$ dieser Gleichung. Irgend eine derselben liefert wegen der Relation
\[
\Delta+k_1f=\varphi^2
\]
die quadratische Form $\varphi=\sqrt{\Delta+k_1f}$; man setzt dieselbe gleich null und bestimmt ihre linearen Factoren $r_x$ und $s_x$, mittelst welcher $f$ in die Normalform (A) transformirt wird. Dieselbe enthält nur mehr gerade Potenzen und kann also durch Quadratwurzeln gelöst werden.

180. \textit{Berechnung des Parameters $m$ der Normalform aus der allgemeinen Form $f=a_x^4$.} Die Herstellung der Normalform von $f$ kann auch noch in anderer Weise erreicht werden, indem wir direct den Parameter $m$ berechnen. Man erhält denselben aus einer Gleichung dritten Grades, die man als eine zweite cubische Resolvente von $f=a_x^4$ auffassen kann. Sind nämlich $i'$ und $j'$ die aus $f=a_x^4$ direct gebildeten Invarianten der biquadratischen Form, so bestehen zwischen ihnen und den Invarianten
\weq{1}{i=2(1+3m^2),\quad j=6m(1-m^2)}
der transformirten Form
\[
f=x^4+6mx^2y^2+y^4
\]
die Relationen:
\weq{2}{\Delta^4i'=i\quad\text{und}\quad\Delta^6j'=j,}
wenn $\Delta$ der Modul der Transformation ist, durch welche die Form $f=a_x^4$ übergeht in $f=x^4+6mx^2y^2+y^4$.


% ===== pdf idx 215, printed 201.0 =====
